\chapter{Overview: GSAS-II Capabilities}
The name GSAS-II derives from the name for the original GSAS program, which is an acronym for General Structure and Analysis System. As such, the original GSAS package was a \textit{system} of many discrete programs that each offered different capabilities. In GSAS-II, we have brought pretty much all of these capabilities into a single GUI. Indeed, as Bob and I have found ourselves in need of some type of computation or visualization, we have added new capabilities into the GSAS-II system. 

In previous chapters we have seen some of these capabilities offered by GSAS-II introduced, but there are are many more and some of these will be described in this section of this book. 
